\documentclass[11pt,a4paper]{article}
\usepackage[margin=2.5cm]{geometry}
\usepackage{amsmath,amssymb,amsthm,mathtools}
\usepackage{booktabs}
\usepackage[colorlinks=true,linkcolor=black,citecolor=black,urlcolor=blue]{hyperref}

\theoremstyle{plain}
\newtheorem{theorem}{Theorem}[section]
\newtheorem{lemma}[theorem]{Lemma}
\newtheorem{proposition}[theorem]{Proposition}
\newtheorem{corollary}[theorem]{Corollary}
\theoremstyle{definition}
\newtheorem{definition}[theorem]{Definition}
\newtheorem{conjecture}[theorem]{Conjecture}
\theoremstyle{remark}
\newtheorem{remark}[theorem]{Remark}

\newcommand{\Tset}{\mathcal{T}}
\newcommand{\Z}{\mathbb{Z}}
\newcommand{\N}{\mathbb{N}}
\newcommand{\Rbar}{\overline{R}}
\newcommand{\Rinf}{R_{\infty}}
\newcommand{\Rmin}{R_{\min}}
\DeclareMathOperator{\Sing}{\mathfrak S}

\title{\textbf{Local constants of the twin-prime propagation map:\\
explicit formulas, mean values, and an unconditional upper bound}}
\author{Dacomb Bierton\\[4pt]
{\normalsize ORCID: 0009-0007-7507-1398}}
\date{September 2026 --- draft}

\begin{document}
\maketitle

\begin{abstract}
Let $\Tset$ be the set of lower twin primes and, for consecutive $p_n<p_{n+1}$
in $\Tset$, let $C_n=p_n+p_{n+1}+1$.  The pair \emph{propagates} if
$C_n\in\Tset$.  The Hardy--Littlewood prediction for the number of
propagating pairs involves, for each gap $g=p_{n+1}-p_n$, the ratio
$R(g)=\Sing_6(g)/\Sing_4(g)$ of the singular series of the six linear forms
$n,\,n+2,\,n+g,\,n+g+2,\,2n+g+1,\,2n+g+3$ to that of the first four.
We prove: (i)~an explicit formula for $R(g)$ as a universal constant
$\Rinf=5.878239\ldots$ times a finite product over the primes dividing
$g(g^2-1)(g^2-4)(g^2-9)$, with the sharp bounds
$\Rmin=2.083577\ldots<R(g)\ll(\log\log g)^2$; (ii)~the mean values
$\frac6G\sum_{g\le G,\,6\mid g}\Sing_4(g)\to 24C_2^2$ and
$\frac6G\sum_{g\le G,\,6\mid g}\Sing_6(g)\to 24C_2^2\,\Rbar$ with
\[
\Rbar \;=\; 9\prod_{q\ge 5}\frac{q^2(q^2-6q+12)}{(q-1)^2(q-2)^2}
\;=\;11.184896\ldots,
\]
which identifies the constant $\kappa=2C_2\Rbar=14.767683\ldots$ in the
quantitative propagation conjecture $N(X)\sim\kappa\int_2^X
dt/\bigl((\log t)^2(\log 2t)^2\bigr)$; (iii)~the unconditional bound
$N(X)\ll X/(\log X)^3$ for the number of propagating pairs with
$p_{n+1}\le X$, one logarithm short of the conjectured order; and (iv)~two
exact identities between singular series: bi-twin chains
$(n,n+2,2n+1,2n+3)$ have the same singular series as prime quadruplets, and
the gap-$6$ propagating configuration has exactly three times the singular
series of a prime sextuplet.  All finite computations are machine-checked
and the numerical predictions are compared with data to $10^{10}$, with $14{,}000$ random twins of up to $150$ digits, and with a proved orbit of the map $t\mapsto2t+d(t)+1$ reaching several hundred digits.
Nothing here proves the twin-prime conjecture; the results are about the
constants of the conjecture and an upper bound in the direction of it.
\end{abstract}

\tableofcontents

\section{Introduction and statement of results}

Let $\Tset=\{3,5,11,17,29,41,59,71,101,\dots\}$ be the set of primes $p$
with $p+2$ prime, enumerated as $p_1<p_2<\cdots$.  In~\cite{BiertonProp}
the author introduced the map $(p_n,p_{n+1})\mapsto C_n:=p_n+p_{n+1}+1$
and conjectured that $C_n\in\Tset$ for infinitely many~$n$
(the \emph{Twin-Prime Propagation Conjecture}); the companion branch
$D_n=C_n+2$ is empty for $n\ge2$~\cite{BiertonThm}.  Write
\[
N(X):=\#\{n\ge1:\ p_{n+1}\le X,\ C_n\in\Tset\}.
\]
Both $p_n$ and $C_n$ are lower twins, so the natural
quantitative version of the conjecture is a Hardy--Littlewood type
asymptotic for $N(X)$.  The purpose of this note is to determine the
constants in that asymptotic exactly, and to prove what can be proved
unconditionally about $N(X)$.

\subsection*{Notation}
For a tuple of linear forms $f_i(n)=a_in+b_i$ ($1\le i\le k$) let
$\nu(q)$ be the number of residues $n\bmod q$ with
$q\mid f_1(n)\cdots f_k(n)$, and let
\begin{equation}\label{eq:sing}
\Sing(f_1,\dots,f_k)=\prod_{q}\Bigl(1-\frac{\nu(q)}q\Bigr)\Bigl(1-\frac1q\Bigr)^{-k}
\end{equation}
be the singular series (the product over primes $q$ converges absolutely
when no form is a constant multiple of another).  The Bateman--Horn
heuristic~\cite{BH} predicts
$\#\{n\le X: \text{all } f_i(n)\text{ prime}\}\sim
\Sing\cdot\int_2^X dt/\prod_i\log f_i(t)$.
We write $C_2=\prod_{q\ge3}q(q-2)/(q-1)^2=0.6601618\ldots$ for the twin
prime constant, so that $\Sing(n,n+2)=2C_2$.

For an even $g$ put
\begin{align*}
\Sing_4(g)&:=\Sing(n,\,n+2,\,n+g,\,n+g+2),\\
\Sing_6(g)&:=\Sing(n,\,n+2,\,n+g,\,n+g+2,\,2n+g+1,\,2n+g+3),\qquad
R(g):=\Sing_6(g)/\Sing_4(g).
\end{align*}
Simultaneous primality of the six forms at $n$ says exactly that
$n,n+g\in\Tset$ and $2n+g+1=C(n,n+g)\in\Tset$.  Both series vanish unless
$6\mid g$ (Lemma~\ref{lem:res}), and for $6\mid g$ the primes $2,3$
contribute the factor $4\cdot\frac94=9$ to $R(g)$ (Lemma~\ref{lem:23}).
Finally let
\[
\Delta(g):=g\,(g^2-1)(g^2-4)(g^2-9)=(g-3)(g-2)(g-1)\,g\,(g+1)(g+2)(g+3).
\]

\subsection*{Results}

\begin{theorem}[Explicit local ratio]\label{thm:explicit}
Let $6\mid g$, $g\ge6$.  For every prime $q\ge5$,
\[
\nu_4(g;q)=\begin{cases}2,& q\mid g,\\ 3,& q\mid g^2-4,\\ 4,&\text{otherwise},\end{cases}
\qquad
\nu_6(g;q)=\begin{cases}\nu_4(g;q),& q\mid (g^2-1)(g^2-9),\\ \nu_4(g;q)+2,&\text{otherwise},\end{cases}
\]
where $\nu_4,\nu_6$ are the root counts of the four- and six-form tuples.
Put $\rho_q(g)=q^2(q-\nu_6)/\bigl((q-1)^2(q-\nu_4)\bigr)$ and
$\gamma_q=q^2(q-6)/\bigl((q-1)^2(q-4)\bigr)$.  Then
\begin{equation}\label{eq:explicit}
R(g)\;=\;\Rinf\;\rho_5(g)\,\rho_7(g)\prod_{\substack{q\ge 11\\ q\mid\Delta(g)}}\frac{\rho_q(g)}{\gamma_q},
\qquad
\Rinf:=9\prod_{q\ge11}\gamma_q=5.87823926\ldots
\end{equation}
In particular $R(g)$ is $\Rinf$ times a finite product determined by
$g\bmod 5$, $g\bmod 7$ and the prime factors $\ge11$ of $\Delta(g)$.
\end{theorem}

\begin{corollary}[Sharp bounds]\label{cor:bounds}
For every $g\equiv0\pmod 6$,
\[
\Rmin:=\Rinf\cdot\frac{25}{48}\cdot\frac{49}{72}=2.08357729\ldots\;<\;R(g)\;\ll\;(\log\log g)^2 .
\]
The constant $\Rmin$ is the infimum of $R(g)$ (approached along
$g\equiv0\pmod5$, $g\equiv\pm2\pmod7$ with all prime factors of $\Delta(g)$
at least $11$ large), and the upper bound is sharp up to the implied
constant: $R(g)\gg(\log\log g)^2$ along $g\equiv 1\pmod q$ for all
$11\le q\le y$, $g\le 6\prod_{q\le y}q$.
\end{corollary}

\begin{theorem}[Mean values over gaps]\label{thm:mean}
As $G\to\infty$,
\begin{align}
\frac6G\sum_{\substack{g\le G\\ 6\mid g}}\Sing_4(g)&\;\longrightarrow\;
\frac{27}{2}\prod_{q\ge5}\frac{q^2(q-2)^2}{(q-1)^4}\;=\;24C_2^2\;=\;10.4595269\ldots,\label{eq:mean4}\\
\frac6G\sum_{\substack{g\le G\\ 6\mid g}}\Sing_6(g)&\;\longrightarrow\;
\frac{243}{2}\prod_{q\ge5}\frac{q^4(q^2-6q+12)}{(q-1)^6}\;=\;24C_2^2\,\Rbar\;=\;116.988726\ldots,\label{eq:mean6}
\end{align}
where
\begin{equation}\label{eq:Rbar}
\Rbar\;=\;9\prod_{q\ge5}\frac{q^2(q^2-6q+12)}{(q-1)^2(q-2)^2}\;=\;11.18489648\ldots
\end{equation}
Consequently the $\Sing_4$-weighted mean of $R(g)$ over gaps tends to
$\Rbar$, and
\begin{equation}\label{eq:kappa}
\kappa:=2C_2\,\Rbar=\frac{27}{2}\prod_{q\ge5}\frac{q^3(q^2-6q+12)}{(q-1)^4(q-2)}=14.76768314\ldots
\end{equation}
\end{theorem}

The constant $\kappa$ is the one that appears in the quantitative form of
the propagation conjecture (Section~\ref{sec:conj}):
\begin{conjecture}[Quantitative propagation]\label{conj:A}
As $X\to\infty$,
\[
N(X)\sim\sum_{p_{n+1}\le X}\frac{R(p_{n+1}-p_n)}{\log C_n\,\log(C_n+2)}
\sim\kappa\int_2^X\frac{dt}{(\log t)^2(\log 2t)^2}.
\]
\end{conjecture}
\noindent
Both forms are compared with the data in Section~\ref{sec:data}: to $10^{10}$
the first predicts $612{,}777$ propagating pairs against $615{,}447$
observed (ratio $1.0044$).  The convergence of the mean of $R(g_n)$ to
$\Rbar$ is slow (the Hardy--Littlewood gap model predicts $11.32$ at height
$10^{10}$; the observed mean is $11.330$), which is why the second form
should be used only asymptotically.

\begin{theorem}[Unconditional upper bound]\label{thm:upper}
$N(X)\ll X/(\log X)^3$.
\end{theorem}
\noindent
The trivial bound is $N(X)\le\#(\Tset\cap[1,X])\ll X/(\log X)^2$ (the
map $n\mapsto C_n$ is injective and $C_n<2X$).  The conjectured order is
$X/(\log X)^4$.

\begin{theorem}[Identities]\label{thm:identities}
\begin{enumerate}
\item[(a)] $\Sing(n,\,n+2,\,2n+1,\,2n+3)=\Sing(n,\,n+2,\,n+6,\,n+8)=4.15118\ldots$:
bi-twin chains of length two are locally exactly as likely as prime
quadruplets.
\item[(b)] $\Sing_6(6)=3\,\Sing(n,\,n+4,\,n+6,\,n+10,\,n+12,\,n+16)=51.8958\ldots$:
a propagating consecutive pair at gap $6$ is locally exactly three times as
likely as a prime sextuplet.
\end{enumerate}
\end{theorem}

\begin{proposition}[Exceptional origins]\label{prop:exceptional}
Call an even $d$ \emph{productive} for $t\in\Tset$ if $t+d\in\Tset$ and
$2t+d+1\in\Tset$~\cite{BiertonGap}.  The tuple
$(d+t,\,d+t+2,\,d+2t+1,\,d+2t+3)$ in the variable $d$ has a fixed prime
divisor if and only if $t=3$; in particular $t=3$ has no productive gap and
every other $t\in\Tset$ has an admissible gap tuple.  If the side condition
``$d-1$ prime'' of~\cite{BiertonGap} is added, the resulting five-form
tuple has a fixed prime divisor exactly for $t\in\{3,5\}$, and for $t=5$ the
only productive $d$ with $d-1$ prime is $d=6$; with the side condition
``$d+1$ prime'' the exceptional set is $\{3\}$.
\end{proposition}

\begin{remark}[What is and is not new]
Theorems~\ref{thm:explicit}--\ref{thm:identities} are statements about
Euler products and residue classes; their proofs are elementary and
complete.  Theorem~\ref{thm:upper} is a routine application of the Selberg
sieve in the form of Halberstam--Richert~\cite[Thm.~5.7]{HR} combined with
Theorem~\ref{thm:mean}; the statement is new because the propagation map
is new, the technique is not.  None of this touches the parity barrier, and
none of it implies the twin-prime conjecture.  The method of
Theorems~\ref{thm:explicit} and~\ref{thm:mean} (root coincidences governed
by cross differences, and averaging of singular series over a shift
parameter) goes back to Hardy--Littlewood and Gallagher~\cite{Gallagher};
the specific tuples, constants and identities do not appear in the
literature as far as the author has been able to determine.
\end{remark}

\section{Residues}

\begin{lemma}\label{lem:res}
If $p\in\Tset$, $p>3$, then $p\equiv5\pmod6$.  If $p<p'$ are in $\Tset$,
both $>3$, then $6\mid p'-p$.  For $g\not\equiv0\pmod6$ the four forms
$n,n+2,n+g,n+g+2$ have a fixed divisor $2$ or $3$, so
$\Sing_4(g)=\Sing_6(g)=0$.
\end{lemma}
\begin{proof}
$p\equiv1\pmod 6$ forces $3\mid p+2>3$; $p\equiv 3\pmod 6$ forces $3\mid p>3$;
even $p$ is not prime.  The difference of two residues $5\bmod6$ is
$0\bmod6$.  For odd $g$ one of $n,n+g$ is even; for
$g\equiv\pm2\pmod6$ the residues $0,-2,-g,-g-2$ cover $\Z/3$.
\end{proof}

\begin{lemma}\label{lem:23}
Let $6\mid g$.  Then $\nu_4(g;2)=\nu_6(g;2)=1$ and $\nu_4(g;3)=\nu_6(g;3)=2$.
Hence the factors at $q=2,3$ of $\Sing_4(g)$, $\Sing_6(g)$ and $R(g)$ are
$8,\ \tfrac{27}{16},\ \tfrac{27}2$ and $32,\ \tfrac{243}{64},\ \tfrac{243}2$
and $4,\ \tfrac94,\ 9$ respectively.
\end{lemma}
\begin{proof}
Modulo $2$: $n$ odd makes all six forms odd (since $g$ is even), so the
only root is $n\equiv0$.  Modulo $3$: with $g\equiv0$ the roots of the four
forms are $0$ and $1$, and the two extra forms $2n+1,2n+3\equiv 2n+1, 2n$
vanish at $n\equiv1$ and $n\equiv0$, adding nothing.  The factors are
$(1-\frac12)2^k=2^{k-1}$ and $(1-\frac23)(\frac32)^k=\frac13(\frac32)^k$
for $k=4,6$.
\end{proof}

\section{The explicit local ratio}

\begin{lemma}[Root coincidences]\label{lem:roots}
Let $6\mid g$ and $q\ge5$ prime.  Modulo $q$ the four forms
$n,n+2,n+g,n+g+2$ have roots $0,-2,-g,-g-2$, and the two extra forms
$2n+g+1,\,2n+g+3$ have roots $r_1=-\tfrac{g+1}2$, $r_2=-\tfrac{g+3}2$.
\begin{enumerate}
\item[(i)] Among $0,-2,-g,-g-2$ there are two coincidences if $q\mid g$,
one if $q\mid g\pm2$, none otherwise; these cases are exclusive.
\item[(ii)] $r_1\neq r_2$.  Moreover
$r_1\in\{0,-2,-g,-g-2\}$ iff $q\mid(g^2-1)(g^2-9)$, and the same holds for
$r_2$ (with the roles of $g\pm1$ and $g\pm3$ interchanged).  Thus either
both extra roots are new or both are old.
\end{enumerate}
\end{lemma}
\begin{proof}
(i) $0\equiv-g$ or $-2\equiv-g-2$ iff $q\mid g$; $0\equiv-g-2$ iff
$q\mid g+2$; $-2\equiv-g$ iff $q\mid g-2$; $0\not\equiv-2$.  A prime
$q\ge5$ divides at most one of $g-2,g,g+2$.

(ii) $r_1\equiv r_2$ iff $q\mid 2$.  Solving $-\frac{g+1}2\equiv a$ for
$a\in\{0,-2,-g,-g-2\}$ gives $q\mid g+1,\ g-3,\ g-1,\ g+3$ respectively;
solving $-\frac{g+3}2\equiv a$ gives $q\mid g+3,\ g-1,\ g-3,\ g+1$.  Both
lists are the divisibility of $q$ into one of $g\pm1,g\pm3$, i.e.\ into
$(g^2-1)(g^2-9)$.  Since $q\ge5$ divides at most one of the four numbers
$g-3,g-1,g+1,g+3$ (their differences are $2,4,6$), the coincidences of $r_1$
and $r_2$ are with two distinct old roots.
\end{proof}

\begin{proof}[Proof of Theorem~\ref{thm:explicit}]
The formulas for $\nu_4$ and $\nu_6$ are Lemma~\ref{lem:roots}.  By
Lemma~\ref{lem:23} and the definition~\eqref{eq:sing},
\[
R(g)=9\prod_{q\ge5}\frac{(1-\nu_6/q)(1-1/q)^{-6}}{(1-\nu_4/q)(1-1/q)^{-4}}
=9\prod_{q\ge5}\frac{q^2(q-\nu_6)}{(q-1)^2(q-\nu_4)}=9\prod_{q\ge5}\rho_q(g).
\]
For $q\ge11$ with $q\nmid\Delta(g)$ we have $(\nu_4,\nu_6)=(4,6)$ and
$\rho_q(g)=\gamma_q$; $\gamma_q=1-9q^{-2}+O(q^{-3})$, so
$\prod_{q\ge11}\gamma_q$ converges and~\eqref{eq:explicit} follows on
dividing out.  (For $q\in\{5,7\}$ the generic case never occurs, since
$g-3,\dots,g+3$ are seven consecutive integers.)
\end{proof}

\begin{proof}[Proof of Corollary~\ref{cor:bounds}]
For fixed $q\ge11$ the four possible values of $\rho_q$ are
\[
\gamma_q=\frac{q^2(q-6)}{(q-1)^2(q-4)},\qquad
\frac{q^2(q-4)}{(q-1)^2(q-2)}\ \ (q\mid g),
\]
\[
\frac{q^2(q-5)}{(q-1)^2(q-3)}\ \ (q\mid g^2-4),\qquad
\frac{q^2}{(q-1)^2}\ \ (q\mid(g^2-1)(g^2-9)),
\]
and $\gamma_q$ is the smallest: $(q-4)^2>(q-2)(q-6)$ and
$(q-5)(q-4)>(q-6)(q-3)$.  Hence every factor in~\eqref{eq:explicit} beyond
$\Rinf\rho_5\rho_7$ is $>1$, and $R(g)\ge\Rinf\min\rho_5\min\rho_7$ with
equality only when no prime $\ge11$ divides $\Delta(g)$.  Enumerating
$g\bmod5$: $\rho_5=\frac{25}{48}$ for $g\equiv0$ and $\frac{25}{16}$
otherwise; $g\bmod7$: $\rho_7=\frac{49}{60},\frac{49}{36},\frac{49}{72}$
for $g\equiv0,\pm1\text{ or }\pm3,\pm2$.  So $R(g)\ge\Rmin$.  Equality
would need $g\equiv0\pmod5$ and every prime factor of $\Delta(g)$ below
$11$; by Sylvester--Schur (a product of $7$ consecutive integers all
$>7$ has a prime factor $>7$) this forces $g\le10$, and $g=6$ has
$g\equiv1\pmod 5$.  So $R(g)>\Rmin$ for all $g$, while choosing $g\equiv0\pmod5$,
$g\equiv\pm2\pmod7$ with all prime factors of $\Delta(g)$ above $11$ large
gives $R(g)\to\Rmin$.

Upper bound: $\rho_q/\gamma_q\le\frac{q^2}{(q-1)^2}\big/\gamma_q=\frac{q-4}{q-6}=1+\frac2{q-6}$,
so $\log R(g)\le\log(\Rinf\rho_5\rho_7)+\sum_{q\mid\Delta(g),\,q\ge11}\frac{2}{q-6}$.
The sum has $\omega(\Delta(g))\le\log\Delta(g)/\log 11\ll\log g$ terms and
is largest when the primes are the smallest possible, i.e.\ at most
$\sum_{11\le q\le y}\frac2{q-6}$ with $\pi(y)\ll\log g$, so
$y\ll\log g\log\log g$ and the sum is $2\log\log y+O(1)\le 2\log\log\log g+O(1)$.
Hence $R(g)\ll(\log\log g)^2$.  Conversely, for $y\ge11$ choose by the
Chinese remainder theorem $g\equiv0\pmod6$ with $g\equiv1\pmod q$ for every
prime $11\le q\le y$ and $0<g\le6\prod_{q\le y}q=e^{(1+o(1))y}$.  Then every
such $q$ divides $g-1\mid\Delta(g)$ with $\rho_q/\gamma_q=1+\frac2{q-6}$,
so $R(g)\ge\Rinf\rho_5\rho_7\prod_{11\le q\le y}(1+\frac2{q-6})\gg(\log y)^2\gg(\log\log g)^2$.
\end{proof}

\section{Mean values}

Throughout this section $k\in\{4,6\}$, $\Sing_k(g)$ is as above and, for a
prime $q\ge5$,
\[
f_q(g):=\Bigl(1-\frac{\nu_k(g;q)}q\Bigr)\Bigl(1-\frac1q\Bigr)^{-k},\qquad u_q(g):=f_q(g)-1,
\]
so that $\Sing_k(g)=c_k\prod_{q\ge5}f_q(g)$ with $c_4=\frac{27}2$, $c_6=\frac{243}2$
(Lemma~\ref{lem:23}).  Note that $f_q(g)$ depends only on $g\bmod q$.

\begin{lemma}[Local means]\label{lem:localmean}
For every prime $q\ge5$,
\[
\frac1q\sum_{a=0}^{q-1}\Bigl(1-\frac{\nu_4(a;q)}q\Bigr)=\frac{(q-2)^2}{q^2},
\qquad
\frac1q\sum_{a=0}^{q-1}\Bigl(1-\frac{\nu_6(a;q)}q\Bigr)=\frac{q^2-6q+12}{q^2}.
\]
\end{lemma}
\begin{proof}
By Lemma~\ref{lem:roots}, $\sum_a(q-\nu_4(a;q))=(q-2)+2(q-3)+(q-3)(q-4)=(q-2)^2$
(one class with $q\mid a$, two with $q\mid a\pm2$, $q-3$ generic).  For
$\nu_6$: the class $a\equiv0$ contributes $q-4$, the classes $a\equiv\pm2$
contribute $2(q-5)$, the four classes $a\equiv\pm1,\pm3$ contribute
$4(q-4)$, and the remaining $q-7$ classes contribute $(q-7)(q-6)$; the total
is $q^2-6q+12$.  This decomposition assumes the seven classes
$0,\pm1,\pm2,\pm3$ are distinct, i.e.\ $q\ge7$; for $q=7$ there are then no
generic classes and the total is $3+4+12=19=49-42+12$.  For $q=5$ the
classes $\pm2$ coincide with $\mp3$ and have $(\nu_4,\nu_6)=(3,3)$, so the
five classes contribute $1+2\cdot2+2\cdot1=7=25-30+12$.  (Both identities
are also verified by enumeration for $5\le q\le 200$ in the accompanying
code.)
\end{proof}

\begin{lemma}[Size of the local deviations]\label{lem:size}
There is an absolute constant $C$ (one may take $C=11$) such that
$|u_q(g)|\le C/q$ for all $q\ge5$, $k\in\{4,6\}$ and all $g$.  If moreover
$q\nmid\Delta(g)$ and $q\ge11$ then $|\log f_q(g)|\le 2k^2/q^2$.
\end{lemma}
\begin{proof}
Write $u_q=\bigl[(1-\nu/q)-(1-1/q)^k\bigr]/(1-1/q)^k$.  For $q\ge7$ the
binomial expansion of $(1-1/q)^k$ is alternating with decreasing terms
(consecutive terms have ratio $\frac{k-j}{(j+1)q}\le k/q<1$), so
$1-k/q\le(1-1/q)^k\le1-k/q+k(k-1)/(2q^2)$ and, since $2\le\nu\le k$, the
numerator lies in $[-k(k-1)/(2q^2),\,(k-\nu)/q]$, of absolute value at most
$4/q$; the denominator is at least $(6/7)^6>0.39$.  The finitely many cases
$q=5$ are checked directly (the largest deviation is $0.53<11/5$).  For the
generic factor,
\[
-\log f_q=-\log(1-k/q)+k\log(1-1/q)=\sum_{j\ge2}\frac{k^j-k}{j\,q^j}
\in\Bigl[0,\ \frac{k^2}{2q^2}\cdot\frac1{1-k/q}\Bigr],
\]
and $1/(1-k/q)\le 11/5$ for $q\ge11$, $k\le6$.
\end{proof}

\begin{lemma}[Averaging a periodic function along multiples of $6$]\label{lem:periodic}
Let $d$ be coprime to $6$, let $u:\Z\to\mathbb R$ be $d$-periodic with
$|u|\le B$ and mean $\bar u=\frac1d\sum_{a=0}^{d-1}u(a)$.  Then for every
$M\ge1$, $\bigl|\sum_{m=1}^{M}u(6m)-M\bar u\bigr|\le 2dB$.  If
$d=q_1\cdots q_r$ is a product of distinct primes and
$u(g)=\prod_i u_{q_i}(g)$ with $u_{q_i}$ $q_i$-periodic, then
$\bar u=\prod_i\bar u_{q_i}$.
\end{lemma}
\begin{proof}
Since $6$ is invertible modulo $d$, $m\mapsto6m$ permutes the residues
modulo $d$, so the sum of $u(6m)$ over any $d$ consecutive $m$ equals
$d\bar u$.  Split $[1,M]$ into $\lfloor M/d\rfloor$ full blocks and a
remainder of fewer than $d$ terms.  The second statement is the Chinese
remainder theorem: $\Z/d\cong\prod\Z/q_i$ and the mean of a product of
functions of the separate coordinates is the product of the means.
\end{proof}

\begin{proof}[Proof of Theorem~\ref{thm:mean}]
Fix $k\in\{4,6\}$, let $G$ be large, $M=\lfloor G/6\rfloor$ and
$Q=\lfloor\log G\rfloor$.  Put $P=\prod_{5\le q\le Q}q$ and
\[
\Sing_k^{(Q)}(g)=c_k\prod_{5\le q\le Q}f_q(g)=c_k\prod_{5\le q\le Q}\bigl(1+u_q(g)\bigr)=c_k\sum_{d\mid P}u_d(g),
\qquad u_d(g):=\prod_{q\mid d}u_q(g).
\]

\emph{Truncation.}  For $6\mid g\le G$,
$\Sing_k(g)/\Sing_k^{(Q)}(g)=\prod_{q>Q}f_q(g)$.  By Lemma~\ref{lem:size}
the primes $q>Q$ not dividing $\Delta(g)$ contribute
$\exp\bigl(O(\sum_{q>Q}q^{-2})\bigr)=1+O(1/Q)$.  The primes $q>Q$ dividing
$\Delta(g)$ number at most $\log|\Delta(g)|/\log Q\le 7\log(G+3)/\log Q$,
and each contributes a factor $1+O(1/q)=1+O(1/Q)$.  Hence, uniformly in
$g\le G$,
\begin{equation}\label{eq:trunc}
\Sing_k(g)=\Sing_k^{(Q)}(g)\Bigl(1+O\Bigl(\frac{\log G}{Q\log Q}\Bigr)\Bigr)
=\Sing_k^{(Q)}(g)\Bigl(1+O\Bigl(\frac1{\log\log G}\Bigr)\Bigr).
\end{equation}

\emph{Averaging.}  Each $u_d$ is $d$-periodic with $|u_d|\le\prod_{q\mid d}C/q=:B_d$
(Lemma~\ref{lem:size}), and $\gcd(d,6)=1$.  By Lemma~\ref{lem:periodic},
\[
\sum_{m=1}^{M}\Sing_k^{(Q)}(6m)=c_k\sum_{d\mid P}\Bigl(M\prod_{q\mid d}\bar u_q+O(dB_d)\Bigr)
=c_kM\prod_{5\le q\le Q}(1+\bar u_q)+O\Bigl(c_k\prod_{5\le q\le Q}(1+C)\Bigr),
\]
because $\sum_{d\mid P}dB_d=\prod_{q\le Q}(1+q\cdot C/q)$.  The error is
$(1+C)^{\pi(Q)}\le(1+C)^{2\log G/\log\log G}=G^{o(1)}=o(M)$.

\emph{Local means.}  By Lemma~\ref{lem:localmean},
$1+\bar u_q=\frac{(q-2)^2}{q^2}\bigl(1-\frac1q\bigr)^{-4}=\frac{q^2(q-2)^2}{(q-1)^4}$
for $k=4$ and
$1+\bar u_q=\frac{q^2-6q+12}{q^2}\bigl(1-\frac1q\bigr)^{-6}=\frac{q^4(q^2-6q+12)}{(q-1)^6}$
for $k=6$; in both cases $\bar u_q=O(q^{-2})$, so
$\prod_{q\ge5}(1+\bar u_q)$ converges and
$\prod_{5\le q\le Q}(1+\bar u_q)=\prod_{q\ge5}(1+\bar u_q)\,(1+O(1/Q))$.

Combining the three steps,
\[
\frac1M\sum_{m\le M}\Sing_k(6m)=c_k\prod_{q\ge5}(1+\bar u_q)\Bigl(1+O\Bigl(\frac1{\log\log G}\Bigr)\Bigr),
\]
which is~\eqref{eq:mean4} and~\eqref{eq:mean6} with the stated Euler
products.  Finally $\prod_{q\ge5}\frac{q(q-2)}{(q-1)^2}=\frac{C_2}{3/4}=\frac{4C_2}3$
gives $\frac{27}2\cdot\frac{16C_2^2}9=24C_2^2$; dividing~\eqref{eq:mean6}
by~\eqref{eq:mean4} gives~\eqref{eq:Rbar}, and $2C_2\Rbar$ with
$2C_2=\frac32\prod_{q\ge5}\frac{q(q-2)}{(q-1)^2}$ gives~\eqref{eq:kappa}.
\end{proof}

\begin{remark}
The error term $O(1/\log\log G)$ is what the crude truncation gives; the
partial means converge much faster in practice (relative error about
$40/G$ at $G\in[6\cdot10^4,1.2\cdot10^5]$).  A Ramanujan-sum
expansion in the style of Montgomery--Soundararajan would give a power
saving, but is not needed here.
\end{remark}

\section{The constant of the propagation conjecture}\label{sec:conj}

For a consecutive pair $(p_n,p_{n+1})$ with gap $g$, simultaneous
membership $p_n,p_{n+1},C_n\in\Tset$ is primality of the six forms at
$n=p_n$.  The Bateman--Horn heuristic for the six forms, divided by that
for the four forms (which describe the event ``$p_n$ and $p_n+g$ are lower
twins'' that we condition on), gives the conditional probability
$R(g)/\bigl(\log C_n\log(C_n+2)\bigr)$ that the pair propagates.  Summing
over pairs gives the first form of Conjecture~\ref{conj:A}.  (Being
consecutive is an additional condition on the pair; heuristically it is
independent of the primality of $C_n,C_n+2$, which live near $2p_n$, and it
is imposed in both numerator and denominator of the conditional
probability.)

For the second form one needs the average of $R(g_n)$ over consecutive
pairs.  In the Hardy--Littlewood model the gap $g$ following a lower twin
near height $x$ has probability proportional to
$\Sing_4(g)\,e^{-\lambda g}$ with $\lambda=2C_2/(\log x)^2$ the density of
$\Tset$ (the factor $\Sing_4(g)$ is the local correlation of the two twins,
the exponential is the probability of no twin in between).  Since
$\lambda\to0$, the average over $g$ of $R(g)$ against this weight tends,
by Theorem~\ref{thm:mean} and partial summation, to the $\Sing_4$-weighted
mean $\Rbar$.  Together with $\#(\Tset\cap[1,x])\sim2C_2x/(\log x)^2$ and
$\log C_n\sim\log 2p_n$ this gives $N(X)\sim2C_2\Rbar\int dt/((\log t)^2(\log2t)^2)$,
i.e.\ $\kappa=2C_2\Rbar$.

The finite-height average
\[
\Rbar(x):=\frac{\sum_{6\mid g}\Sing_6(g)e^{-\lambda g}}{\sum_{6\mid g}\Sing_4(g)e^{-\lambda g}},
\qquad\lambda=\frac{2C_2}{(\log x)^2},
\]
converges to $\Rbar$ only like the partial means of Theorem~\ref{thm:mean}
at $G\approx1/\lambda$.  Table~\ref{tab:Rx} lists it; the observed mean of
$R(g_n)$ over all consecutive pairs with $p_n\le10^{10}$ is $11.3297$
($11.3467$ to $10^9$), and over the dyadic blocks $[2^{28},2^{29})$ and
$[2^{33},2^{34})$ it is $11.335$ and $11.320$, in agreement with the model
at those heights and drifting in the predicted direction.

\begin{table}[h]\centering
\begin{tabular}{@{}rrr@{}}\toprule
$x$ & $\Rbar(x)$ & mean gap $1/\lambda$-weighted\\ \midrule
$10^{6}$ & $11.448$ & $154$\\
$10^{8}$ & $11.366$ & $268$\\
$10^{9}$ & $11.340$ & $337$\\
$10^{10}$ & $11.319$ & $414$\\
$10^{12}$ & $11.290$ & $592$\\
$10^{20}$ & $11.238$ & $1623$\\
$10^{50}$ & $11.201$ & $\approx10^4$\\
$\infty$ & $11.18490$ & \\ \bottomrule
\end{tabular}
\caption{Model mean of $R(g_n)$ at height $x$ and its limit $\Rbar$.}\label{tab:Rx}
\end{table}

\section{The unconditional upper bound}

We use the Selberg upper-bound sieve in the following form.

\begin{theorem}[Halberstam--Richert {\cite[Thm.~5.7]{HR}}]\label{thm:HR}
Let $F_i(n)=a_in+b_i$ ($1\le i\le k$) be linear forms with integer
coefficients, $a_i\ge1$, such that
$E:=\prod_ia_i\prod_{r<s}(a_rb_s-a_sb_r)\neq0$, and let $\rho(q)$ be the
number of roots of $\prod_iF_i$ modulo $q$.  Suppose $\rho(q)<q$ for all
primes $q$.  Then for $1<y\le z$,
\[
\#\{z-y<n\le z:\ F_1(n),\dots,F_k(n)\text{ all prime}\}
\le 2^kk!\,\prod_q\Bigl(1-\frac{\rho(q)}q\Bigr)\Bigl(1-\frac1q\Bigr)^{-k}\frac{y}{(\log y)^k}
\bigl(1+O(\varepsilon)\bigr),
\]
where $\varepsilon=(\log\log3y+\log\log3|E|)/\log y$ and the implied
constant depends only on $k$.
\end{theorem}

\begin{lemma}\label{lem:adm6}
For every $g\equiv0\pmod6$ the six forms $n,n+2,n+g,n+g+2,2n+g+1,2n+g+3$
satisfy $\nu_6(g;q)<q$ for every prime $q$.
\end{lemma}
\begin{proof}
$\nu_6=1<2$ and $\nu_6=2<3$ by Lemma~\ref{lem:23}; for $q\ge5$,
$\nu_6\le\nu_4+2\le6$, which is $<q$ for $q\ge7$.  For $q=5$ use
Lemma~\ref{lem:roots}: $5$ divides exactly one of $g-2,g-1,g,g+1,g+2$.
If $5\mid g\pm1$ then $\nu_6=\nu_4\le4$.  If $5\mid g\pm2$ then also
$5\mid g\mp3$ (as $g+2\equiv g-3$), so again $\nu_6=\nu_4=3$.  If $5\mid g$
then $\nu_4=2$ and $\nu_6\le4$.  In every case $\nu_6(g;5)\le4<5$.
\end{proof}

\begin{proof}[Proof of Theorem~\ref{thm:upper}]
Let $X$ be large and $G=\lfloor(\log X)^3\rfloor$.  Every $n$ counted by
$N(X)$ has $p_{n+1}\le X$; write $g_n=p_{n+1}-p_n$.

\emph{Large gaps.}  The gaps $g_n$ with $p_{n+1}\le X$ are disjoint
intervals inside $[3,X]$, so $\sum_{p_{n+1}\le X}g_n\le X$ and
\[
\#\{n:\ p_{n+1}\le X,\ g_n>G\}\le X/G\ll X/(\log X)^3 .
\]

\emph{Small gaps.}  Fix $g\le G$ with $6\mid g$ (other $g$ do not occur by
Lemma~\ref{lem:res}, apart from the single pair $(3,5)$).  If $g_n=g$ and
the pair propagates, then $n':=p_n\le X$ makes the six forms
$n',n'+2,n'+g,n'+g+2,2n'+g+1,2n'+g+3$ simultaneously prime.  Apply
Theorem~\ref{thm:HR} with $k=6$, $y=z=X$: the hypothesis $\rho(q)<q$ is
Lemma~\ref{lem:adm6}; the coefficients are $a_i\in\{1,2\}$ and
$|b_i|\le g+3$, so $|E|\le2^6(2(g+3))^{15}\le(\log X)^{60}$ and
$\log\log3|E|\ll\log\log\log X$.  Hence
\[
\#\{n:\ p_{n+1}\le X,\ g_n=g,\ C_n\in\Tset\}\ll\Sing_6(g)\frac{X}{(\log X)^6},
\]
uniformly in $g\le G$.  Summing over $6\mid g\le G$ and using
Theorem~\ref{thm:mean} (which gives $\sum_{6\mid g\le G}\Sing_6(g)\ll G$),
\[
\#\{n:\ p_{n+1}\le X,\ g_n\le G,\ C_n\in\Tset\}\ll\frac{X}{(\log X)^6}\sum_{\substack{g\le G\\6\mid g}}\Sing_6(g)
\ll\frac{GX}{(\log X)^6}\ll\frac{X}{(\log X)^3}.
\]
Adding the two parts proves the theorem.
\end{proof}

\begin{remark}
The loss of one logarithm against the conjectured $X/(\log X)^4$ comes
entirely from the trivial treatment of large gaps: one would need
$\#\{p_n\le X:\ g_n>(\log X)^{3}\}\ll X/(\log X)^{4}$, a statement about
twin primes in short intervals that is not available.  Any improvement of
the large-gap count improves Theorem~\ref{thm:upper} correspondingly.
\end{remark}

\section{Identities and exceptional origins}

\begin{proof}[Proof of Theorem~\ref{thm:identities}]
(a) For the bi-twin tuple $n,n+2,2n+1,2n+3$ the roots modulo $q\ge5$ are
$0,-2,-\frac12,-\frac32$; their pairwise differences are
$2,\frac12,\frac32,\frac32,\frac12,1$ (up to sign), none of which vanishes
modulo $q\ge5$, so $\nu(q)=4$.  For the quadruplet $n,n+2,n+6,n+8$ the
roots $0,-2,-6,-8$ have differences $2,6,8,4,6,2$, again nonzero modulo
$q\ge5$, so $\nu(q)=4$.  Modulo $2$ both tuples have $\nu=1$; modulo $3$
both have $\nu=2$ (roots $0,1$ in each case).  The Euler products
coincide factor by factor.

(b) The propagation tuple with $g=6$ has, by Theorem~\ref{thm:explicit},
$\nu_6(6;5)=4$ (as $5\mid g-1$), $\nu_6(6;7)=4$ (as $7\mid g+1$), and
$\nu_6(6;q)=6$ for $q\ge11$ because $\Delta(6)=2^{6}\cdot3^{4}\cdot5\cdot7$
has no prime factor $\ge11$; and $\nu=1,2$ at $q=2,3$.  The sextuplet
$n,n+4,n+6,n+10,n+12,n+16$ has roots $0,-4,-6,-10,-12,-16$; modulo $5$
these are $\{0,1,4,0,3,4\}$, four classes; modulo $7$ they are
$\{0,3,1,4,2,5\}$, six classes; the pairwise differences
$4,6,10,12,16,2,6,8,12,4,6,10,2,6,4$ are not divisible by any prime
$\ge11$, so $\nu=6$ for $q\ge11$; and $\nu=1,2$ at $q=2,3$.  The two
root-count profiles agree except at $q=7$, where the local factors are
$(1-\frac47)(\frac67)^{-6}$ and $(1-\frac67)(\frac67)^{-6}$, whose ratio is
$3$.
\end{proof}

\begin{proof}[Proof of Proposition~\ref{prop:exceptional}]
Modulo $2$ all four forms $d+t,d+t+2,d+2t+1,d+2t+3$ are odd for even $d$.
Modulo $3$, for $t\equiv2$ (every $t\in\Tset$, $t>3$) the roots in $d$ are
$-t\equiv1$, $-t-2\equiv2$, $-2t-1\equiv1$, $-2t-3\equiv2$, so $d\equiv0$
survives; for $t=3$ the roots are $0,1,2,0$, covering $\Z/3$.  For $q\ge5$
there are at most four roots.  Hence the tuple has a fixed divisor iff
$t=3$; and for $t=3$ any $d$ with $3+d\in\Tset$ has $d\equiv2\pmod6$, so
$2\cdot3+d+1\equiv0\pmod 3$ and exceeds $3$: no productive gap exists.

With the fifth form $d-1$: modulo $5$ and $t=5$ the roots are
$0,3,4,2,1$, all of $\Z/5$, so exactly one of $d+5,d+7,d+11,d+13,d-1$ is
divisible by $5$ for every $d$; primality forces that one to equal $5$,
which happens only for $d-1=5$, i.e.\ $d=6$ (and indeed
$11,13,17,19$ are prime).  For $t\in\Tset$, $t\ge11$, one has
$t\bmod30\in\{11,17,29\}$ and a direct enumeration gives $\nu_5\le4$ in
all three cases and for both side conditions; modulo $3$ the fifth form
$d\mp1\equiv\mp1$ adds no root at $d\equiv0$; for $q\ge7$ there are at most
five roots.  With the form $d+1$ and $t=5$ the roots modulo $5$ are
$0,3,4,2,4$, so the tuple is admissible.  This gives the exceptional sets
$\{3,5\}$ and $\{3\}$.
\end{proof}

\section{Numerical verification}\label{sec:data}

All constants were evaluated as Euler products truncated at
$Q=10^8$ with the analytic tail $\sum_{q>Q}q^{-2}\approx(1-1/\log Q)/(Q\log Q)$
added; changing $Q$ from $10^7$ to $10^8$ moves the values by less than
$2\cdot10^{-9}$.  As a check, the product form of $24C_2^2$ in~\eqref{eq:mean4}
agrees with the value computed from the known digits of $C_2$ to
$10^{-10}$.

\begin{center}
\begin{tabular}{@{}lr@{}}\toprule
$\Rinf=9\prod_{q\ge11}\gamma_q$ & $5.87823926$\\
$\Rmin=\Rinf\cdot\frac{25}{48}\cdot\frac{49}{72}$ & $2.08357729$\\
$R(6)=\Rinf\cdot\frac{25}{16}\cdot\frac{49}{36}$ & $12.5014637$\\
$24C_2^2$ & $10.4595270$\\
$\Rbar$ & $11.1848965$\\
$24C_2^2\Rbar$ & $116.988726$\\
$\kappa=2C_2\Rbar$ & $14.7676831$\\
$\Sing(\text{bi-twin})=\Sing(\text{quadruplet})$ & $4.15118075$\\
$\Sing_6(6)=3\Sing(\text{sextuplet})$ & $51.8958333$\\ \bottomrule
\end{tabular}
\end{center}

Theorem~\ref{thm:explicit} was checked against the direct product over
the same primes for $g$ up to $9{,}699{,}690$ (relative difference $<10^{-12}$),
and the root-count table against enumeration for all $6\mid g\le3000$,
$5\le q\le300$.  The minimum of $R(g)$ over $6\mid g\le2\cdot10^5$ is
$2.083876$ at $g=195{,}540$ (so $R(g)>\Rmin$ with the infimum approached),
and the maximum is $42.51$ at $g=194{,}466$.  Lemma~\ref{lem:localmean}
was checked by enumeration for $5\le q\le200$, and the partial means of
Theorem~\ref{thm:mean} to $G=1.2\cdot10^5$ are $10.4560$ and $116.970$
against the limits $10.4595$ and $116.989$.

\subsection*{The propagation conjecture to $10^{10}$}

Among the $27{,}412{,}679$ lower twins up to $10^{10}$ there are
$615{,}447$ propagating consecutive pairs; the first form of
Conjecture~\ref{conj:A} predicts $612{,}777.0$ (ratio $1.0044$) and the
$\kappa$-integral $612{,}694.8$ (ratio $1.0045$).  The per-dyadic-block
ratios are within $\pm1.7\%$ from $2^{22}$ on and within $\pm0.7\%$ from
$2^{28}$ on.  The only pair with $D_n\in\Tset$ is $(3,5)$, as
in~\cite{BiertonThm}.  Restricting to gap $6$: there
are $180{,}529$ consecutive pairs at gap $6$ with $p_n\le10^{10}$ (prime
quadruplets) against the Hardy--Littlewood prediction $181{,}065.1$, and
$4{,}527$ of them propagate against the Bateman--Horn prediction
$\Sing_6(6)\int_2^{10^{10}}dt/\bigl((\log t)^2(\log(t+6))^2(\log(2t+7))^2\bigr)=4{,}540.6$
for the six forms (ratio $0.997$), with $\Sing_6(6)=3\,\Sing(\text{sextuplet})$
from Theorem~\ref{thm:identities}.  Per gap, the ratios observed/predicted
for the twenty most frequent gaps $g\le120$ lie between $0.945$ and $1.036$.

\subsection*{The minimal productive gap}

For $t\in\Tset$, $t\ge5$, let $d(t)$ be the least $d\equiv0\pmod6$ that is
productive for $t$ in the sense of~\cite{BiertonGap} (i.e.\ $t+d\in\Tset$,
$2t+d+1\in\Tset$, and $d-1$ or $d+1$ prime).  The note~\cite{BiertonGap}
conjectures that $d(t)$ exists for every $t\ge5$; the Bateman--Horn
heuristic for the five-form tuple in the variable $d$, with the fixed
divisors of Proposition~\ref{prop:exceptional} removed, gives an expected
number $E_t(Y)$ of productive gaps $\le Y$ of order $Y/(\log t)^4$, and the
Poisson model $\Pr(d(t)>Y)=e^{-E_t(Y)}$ then predicts
$d(t)\asymp(\log t)^4\log\log t$ on average and
$d(t)\le(\log t)^5\log\log t$ for every $t\ge5$ (the refined form of the
conjecture used in the repository).  Three independent data sets test this.

\emph{Exhaustive range.}  Every one of the $27{,}412{,}678$ lower twins
$5\le t\le10^{10}$ has a productive gap; the largest is
$d(9{,}432{,}106{,}511)=1{,}728{,}906$, and the largest value of
$d(t)/((\log t)^5\log\log t)$ is $0.144$, at $t=324{,}301{,}421$.  The mean
of $d(t)/((\log t)^4\log\log t)$ is $0.0573$ or $0.0574$ in every dyadic
block from $[2^{26},2^{27})$ to $[2^{33},2^{34})$, and the model mean
$\sum_Y e^{-E_t(Y)}$ matches the observed block means to within $10\%$
throughout and to $0.1\%$ in the last block.  The side condition is met by
$d-1$ alone for $43.3\%$ of the twins, by $d+1$ alone for $42.9\%$, and by
both for $13.8\%$.

\emph{Independent samples at large height.}  For each of the sizes
$20,30,40,60,80,100,150$ decimal digits, $2{,}000$ lower twins were drawn
uniformly at random (random start in the decade, forward scan) and $d(t)$
was computed exactly by sieving the candidates $d=6m$ with all primes up to
$10^5$ before testing survivors.  Under the Poisson model the values
$u=e^{-E_t(d(t))}$ are uniform on $(0,1)$.  Over the $14{,}000$ samples the
mean of $u$ is $0.5033$ and the Kolmogorov--Smirnov distance to the uniform
distribution is $0.0086$, against a $5\%$ critical value of $0.0115$; the
seven per-size distances are $0.018,0.039,0.023,0.019,0.028,0.017,0.020$
(critical value $0.030$), and the ratios
$\sum d(t)/\sum\mathbb E[d(t)]$ lie between $0.954$ and $1.022$.  The
normalised mean $d(t)/((\log t)^4\log\log t)$ is $0.056$--$0.061$ at every
size, the same value as in the exhaustive range.

\emph{The constructive orbit.}  Iterating $t\mapsto2t+d(t)+1$ from $t=5$
produces a single orbit through $\Tset$ (the ``twin-gap chain''
of~\cite{BiertonGap}).  It was followed for 1{,}134 steps, to a $t$ of
343 decimal digits.  Every prime along it ($t$, $t+2$, $t+d$,
$t+d+2$ at each step, and the final pair) was proved prime with the APR--CL
test of PARI/GP, and ECPP certificates for the final pair are distributed
with the repository, so the orbit does not rest on probable-prime tests.
Along the orbit the values $u=e^{-E_t(d(t))}$ have mean 0.516 and
Kolmogorov--Smirnov distance 0.038 to the uniform distribution ($5\%$
critical value 0.040); the largest value of
$d(t)/((\log t)^5\log\log t)$ is $0.556$, attained at the first step $t=5$,
and it is below 0.005 for every $t>10^{12}$.  The orbit therefore
behaves, as far as $d(t)$ is concerned, like a sequence of independent
twins of the same sizes, across some 340 orders of magnitude.

All checks are reproduced in the accompanying repository%
\footnote{\url{https://github.com/dacomb-bierton/Local-Constants-Twin-Prime-Propogation}} by the commands
\texttt{certify}, \texttt{theorems}, \texttt{propagation} and \texttt{gaps}
of the \texttt{twinconj} package and by the scripts \texttt{extend\_chain},
\texttt{sample\_gaps} and \texttt{certify\_chain}; the two exhaustive runs
to $10^{10}$ need about $40$\,GB of memory.

\section{Status}

\begin{center}
\begin{tabular}{@{}p{0.58\textwidth}p{0.36\textwidth}@{}}\toprule
\textbf{Claim} & \textbf{Status}\\ \midrule
Explicit formula and sharp bounds for $R(g)$ & Theorem~\ref{thm:explicit}, Corollary~\ref{cor:bounds}: proved.\\
Mean values of $\Sing_4,\Sing_6$ over gaps; closed forms for $\Rbar$, $\kappa$ & Theorem~\ref{thm:mean}: proved.\\
$N(X)\ll X/(\log X)^3$ & Theorem~\ref{thm:upper}: proved (Selberg sieve).\\
Bi-twin $=$ quadruplet, $\Sing_6(6)=3\,\Sing(\text{sextuplet})$ & Theorem~\ref{thm:identities}: proved.\\
Exceptional origins $\{3\}$, $\{3,5\}$ & Proposition~\ref{prop:exceptional}: proved.\\
$N(X)\to\infty$ (propagation conjecture) & Open; follows from Hypothesis~H~\cite{BiertonThm}.\\
$N(X)\sim\kappa\int dt/((\log t)^2(\log2t)^2)$ & Conjecture~\ref{conj:A}; consistent with data to $10^{10}$ (ratio $1.0044$).\\
Twin-prime conjecture & Open; not addressed by anything here.\\ \bottomrule
\end{tabular}
\end{center}

\begin{thebibliography}{99}
\bibitem{BH} P.~T. Bateman and R.~A. Horn, A heuristic asymptotic formula
concerning the distribution of prime numbers, \textit{Math.\ Comp.} 16
(1962), 363--367.
\bibitem{BiertonProp} D.~Bierton, A twin-prime propagation conjecture,
Zenodo, August 2026, \url{https://doi.org/10.5281/zenodo.22017371}.
\bibitem{BiertonGap} D.~Bierton, A twin-gap existence conjecture, Zenodo,
August 2026, \url{https://doi.org/10.5281/zenodo.22058355}.
\bibitem{BiertonH} D.~Bierton, Structural reduction of two twin-prime
conjectures and conditional resolution under Hypothesis~H, Zenodo, August
2026, \url{https://doi.org/10.5281/zenodo.22063097}.
\bibitem{BiertonThm} D.~Bierton, Theorems on twin-prime propagation and a
redefinition of two conjectures, Zenodo, 31 August 2026,
\url{https://doi.org/10.5281/zenodo.22206095}.
\bibitem{Gallagher} P.~X. Gallagher, On the distribution of primes in
short intervals, \textit{Mathematika} 23 (1976), 4--9.
\bibitem{HR} H.~Halberstam and H.-E. Richert, \textit{Sieve Methods},
Academic Press, London, 1974.
\bibitem{HL} G.~H. Hardy and J.~E. Littlewood, Some problems of `Partitio
Numerorum'. III, \textit{Acta Math.} 44 (1923), 1--70.
\bibitem{MS} H.~L. Montgomery and K.~Soundararajan, Primes in short
intervals, \textit{Comm.\ Math.\ Phys.} 252 (2004), 589--617.
\bibitem{OEIS} OEIS Foundation Inc., Sequence A066388 (bi-twin chains),
\url{https://oeis.org/A066388}.
\end{thebibliography}

\end{document}
